\documentclass{article} 
\input{../../../lib.sty} 
 
\title{\texttt{fn sample\_bernoulli\_rational}} 
\author{\MichaelShoemateAuthor} 
 
\begin{document} 
\maketitle 
 
This document proves that the implementations of \rustdoc{traits/samplers/bernoulli/fn}{sample\_bernoulli\_rational} in \asOfCommit{mod.rs}{f5bb719}  
satisfies its proof definition. 
 
\section{Hoare Triple} 
\subsection*{Precondition} 
\subsubsection*{Compiler-Verified} 
\begin{itemize} 
    \item Variable \texttt{prob} must be of type \texttt{RBig}, a rational of bignums. 
\end{itemize} 
 
\subsubsection*{Caller-Verified} 
None 
 
\subsubsection*{Pseudocode} 
 
\lstinputlisting[language=Python,firstline=2,escapechar=|]{./pseudocode/sample_bernoulli_rational.py} 
 
\subsubsection*{Postcondition} 
\begin{theorem} 
    \label{sample-bernoulli} 
 
    For any setting of the input parameters, 
    returns an error if there is a lack of system entropy or \texttt{prob} is not in $[0, 1]$, 
    returns true with probability \texttt{prob}, otherwise returns false. 
\end{theorem} 
 
\begin{proof}  
At \ref{prob-parts}, \texttt{into\_parts} returns the signed numerator and unsigned nonzero denominator of \texttt{prob}. 
At \ref{numer-parts}, \texttt{into\_parts} returns the sign and magnitude of \texttt{numer}. 
 
An error is raised if the sign is negative (if \texttt{prob} is negative)  
or if the numerator is greater than the denominator (if \texttt{prob} is greater than one). 
 
Since \texttt{RBig.into\_parts} never returns a zero denominator (this would be an invalid rational), 
the precondition of \texttt{sample\_uniform\_ubig\_below} is met, 
therefore by the postcondition the return value is a uniform sample in $[0, denom)$. 
 
By countable additivity, the probability that the uniform sample is less than the numerator is exactly \texttt{prob}, 
as there are \texttt{numer} possible disjoint outcomes each with probability $1 / \texttt{denom}$. 
It has been shown that the function returns true with probability \texttt{prob}, otherwise returns false. 
\end{proof} 
 
\end{document} 
